\section{Analysis}
\label{sec:analysis}

\subsection{Why Does A3 Work?}
\label{sec:why_a3_works}

The primary conceptual challenge in post-training low-rank factorization is that unweighted truncated singular value decomposition minimizes the unweighted Frobenius error $\|W - \tilde{W}\|_F^2 = \sum_{i,j} (W_{ij} - \tilde{W}_{ij})^2$. By treating all matrix coordinates and input dimensions identically, standard SVD implicitly presumes an isotropic input activation distribution: $x \sim \mathcal{N}(0, I)$. However, representations inside modern multi-head transformer networks violate this assumption by orders of magnitude. The efficacy of A3 arises from the coordinated operation of three complementary mechanisms: empirical Fisher sensitivity reweighting, adaptive per-layer rank escalation with deterministic fallback, and analytical LayerNorm recalibration. Below, we conduct an extensive analytical and empirical dissection of each constituent component.

\subsubsection{Fisher Weighting Distorts the Singular Value Spectrum to Protect Salient Subspaces}
\label{sec:fisher_weighting_analysis}

In deep contextualized transformer models, activation distributions exhibit extreme coordinate-wise anisotropy. Recent inquiries into post-training quantization, such as AWQ (Lin et al., 2023) and Wanda (Sun et al., 2023), established that a minuscule subset of hidden channels acts as persistent feature outliers, displaying variance orders of magnitude greater than the median coordinate. In standard low-rank approximation, truncation errors allocated along these high-variance dimensions cause disproportionate damage to the network's downstream representations because output perturbations scale directly with input channel energy: $\mathbb{E}[\|(W - \tilde{W})x\|_2^2] = \mathrm{Tr}((W - \tilde{W}) C_X (W - \tilde{W})^\top)$.

To address this distortion without resorting to expensive non-diagonal metric inversions, A3 utilizes the empirical Fisher diagonal $F_i = \frac{1}{N}\sum_{t=1}^N (x_{t,i})^2$ as a computationally tractable proxy for parameter importance. In a properly trained network, coordinates exhibiting high average squared magnitude correspond directly to dominant gradients under standard quadratic loss bounds. Figure~\ref{fig:fisher_diagonal} plots the sorted empirical Fisher diagonal energy distribution for an intermediate attention projection in DistilBERT ($d=768$). The empirical distribution is profoundly heavy-tailed: the top 10\% of feature coordinates concentrate approximately 85\% of the total sensitivity norm, whereas the lowest 50\% of coordinates account for less than 2\% of aggregate sensitivity.

\begin{figure}[H]
\centering
\includegraphics[width=\columnwidth]{figures/fig_8_1.pdf}
\caption{Sorted empirical Fisher diagonal energy distribution for an intermediate DistilBERT attention layer. The top 10\% of feature coordinates concentrate $\sim$85\% of the total sensitivity norm, justifying activation-weighted factorization. Source: \texttt{results/a3\_tests/fixed\_v4/a3\_stats.json}.}
\label{fig:fisher_diagonal}
\end{figure}

Mathematically, multiplying the weight matrix $W$ by the diagonal scaling factor $S = \mathrm{diag}(\mathbf{s})$, where $\mathbf{s} = (\sqrt{F + \epsilon}) / \mathrm{RMS}(\sqrt{F + \epsilon})$, rescales the column vectors of $W$ proportionally to their empirical activation energy. When truncated SVD is computed on the warped matrix $W_\text{weighted} = W S = U \Sigma V^\top$, the Eckart-Young-Mirsky theorem guarantees that the optimal low-rank projection minimizes the weighted reconstruction objective $\|(W - \tilde{W})S\|_F^2$. Consequently, coordinates possessing large Fisher values are magnified during the singular value decomposition, forcing the principal left and right singular vectors to span the subspace aligned with these high-energy coordinates. Inverting the scaling factor via $\tilde{W} = (U_r \Sigma_r V_r^\top) S^{-1}$ restores the unscaled coordinate system. Crucially, reconstruction errors are concentrated almost exclusively in coordinates where $S_{ii}$ is negligible, ensuring that the expected output perturbation $\mathbb{E}[\|\Delta W x\|^2]$ is minimized. This spectral reshaping directly explains why Fisher weighting retains critical classification boundaries where vanilla SVD collapses.

\subsubsection{Adaptive Rank Allocation Resolves Layer Heterogeneity and Bottleneck Dynamics}
\label{sec:adaptive_rank_analysis}

A persistent deficiency of historical matrix factorization schemes in deep learning is the enforcement of a static, uniform rank across all network layers. In multi-layer architectures, intrinsic representation dimensionality evolves non-linearly across depth: early layers process localized, syntax-heavy token encodings with relatively high rank, middle layers exhibit subspace clustering, and late layers form task-specific decision manifolds. Forcing an invariant rank across these distinct regimes inevitably produces severe under-parameterization in critical bottleneck projections while over-allocating rank capacity to redundant operators.

\begin{figure}[H]
\centering
\includegraphics[width=\columnwidth]{figures/fig_8_2.pdf}
\caption{Allocated low-rank dimension per transformer layer under adaptive rank escalation subject to $r_{\mathrm{ceiling}} \le 256$ and threshold $\tau = 0.40$. Source: \texttt{results/a3\_tests/fixed\_v4/a3\_stats.json}.}
\label{fig:rank_per_layer}
\end{figure}

A3 overcomes layer heterogeneity by coupling progressive rank escalation with an empirical Frobenius quality gate. As illustrated in Figure~\ref{fig:rank_per_layer}, early query and key projections in DistilBERT tolerate substantial compression, achieving satisfactory representation fidelity at ranks $r = 64$ to $128$. Conversely, intermediate attention projections encounter steep error curves, escalating through $r=192$ up to the ceiling $r_{\mathrm{ceiling}} = 256$. Enforcing $r_{\mathrm{ceiling}} \le 256$ guarantees that the factorized parameter count $r(d_\text{in} + d_\text{out})$ never breaches the break-even dimension $d_\text{in} d_\text{out} / (d_\text{in} + d_\text{out}) = 384$, entirely eliminating negative parameter savings. Whenever a layer fails to achieve relative reconstruction error below $\tau = 0.40$ at $r_{\mathrm{ceiling}}$, the algorithm triggers deterministic fallback, leaving the original uncompressed dense parameters intact. This adaptive bifurcation insulates the network against catastrophic layer collapse.

\subsubsection{Closed-Form LayerNorm Recalibration Absorbs Downstream Variance Truncation}
\label{sec:ln_recalibration_analysis}

When a dense linear transformation $y = W x$ is replaced by its rank-$r$ approximation $\tilde{y} = \tilde{W} x$, the omission of trailing singular values $\sigma_{r+1}, \dots, \sigma_d$ introduces systematic variance shrinkage:
\begin{equation}
\mathbb{E}[\|\tilde{W} x\|_2^2] = \sum_{i=1}^r \sigma_i^2 < \sum_{i=1}^d \sigma_i^2 = \mathbb{E}[\|W x\|_2^2]
\end{equation}
In deep residual architectures, this attenuation compounds across successive layers, systematically shifting the operating regime of subsequent LayerNorm modules. We empirically measure an average reduction in pre-normalization feature standard deviation of up to 15\% immediately following low-rank projection. Because standard LayerNorm normalizes inputs by $\sqrt{\mathrm{Var}(x) + \epsilon}$, a uniform shrinkage in input variance amplifies numerical noise near the epsilon floor and alters the gradient scale during downstream adaptation.

A3 counteracts this distributional shift in closed form without gradient backpropagation. By tracking the empirical token standard deviation before compression ($\sigma_\text{target}$) and immediately following factorization ($\sigma_\text{current}$), we update the affine gain vector according to:
\begin{equation}
\gamma_\text{new} = \gamma_\text{old} \odot \left(\frac{\sigma_\text{current}}{\sigma_\text{target}}\right)
\end{equation}
This affine adjustment directly compensates for the post-factorization signal shrinkage, restoring the dynamic range of activations entering subsequent non-linearities. Empirically, this instantaneous rescaling closes over 60\% of the initial accuracy deficit before any gradient recovery steps are executed.

\subsection{Comparison to Unweighted SVD}
\label{sec:svd_comparison}

To rigorously quantify the isolated benefit of Fisher subspace weighting, we execute an ablation comparing A3 against classical unweighted truncated SVD across identical test architectures at a uniform rank constraint $r=128$.

\begin{table}[H]
\centering
\small
\begin{tabular}{lllll}
\toprule
\textbf{Model} & \textbf{Factorization Scheme} & \textbf{SST-2 Accuracy} & \textbf{Reconstruction Error} & \textbf{F1 Score} \\
\midrule
DistilBERT & Unweighted Truncated SVD & 86.42\% & 0.586 & 0.8681 \\
DistilBERT & \textbf{A3 (Fisher-Weighted)} & \textbf{90.02\%} & \textbf{0.375} & \textbf{0.9043} \\
RoBERTa & Unweighted Truncated SVD & 73.15\% & 0.542 & 0.7410 \\
RoBERTa & \textbf{A3 (Fisher-Weighted)} & \textbf{81.19\%} & \textbf{0.381} & \textbf{0.8322} \\
GPT-2 & Unweighted Truncated SVD & 58.70\% & 0.512 & 0.6514 \\
GPT-2 & \textbf{A3 (Fisher-Weighted)} & \textbf{63.19\%} & \textbf{0.341} & \textbf{0.7216} \\
\bottomrule
\end{tabular}
\caption{Systematic empirical comparison between standard unweighted truncated SVD and Fisher-weighted A3 at fixed rank $r=128$ under identical calibration budgets. Sources: \texttt{results/a3\_tests/fixed\_v4/a3\_stats.json}, \texttt{results/unified\_a3/roberta/roberta-base\_a3\_full\_metrics.json}, \texttt{results/unified\_a3/gpt2/gpt2\_a3\_full\_metrics.json}.}
\label{tab:svd_vs_a3}
\end{table}

The empirical results compiled in Table~\ref{tab:svd_vs_a3} confirm the decisive role of activation weighting. On DistilBERT, Fisher weighting reduces the mean relative Frobenius reconstruction error from 0.586 down to 0.375, delivering a $+3.60\%$ absolute accuracy boost (from 86.42\% to 90.02\%) and elevating the F1 metric from 0.8681 to 0.9043. On RoBERTa-base, the performance divergence is even more pronounced: unweighted SVD induces catastrophic disruption in intermediate representation geometry, causing accuracy to collapse to 73.15\%, whereas A3 sustains 81.19\% ($\Delta +8.04\%$). 

Spectral analysis reveals that RoBERTa activations exhibit greater coordinate kurtosis than DistilBERT due to deeper layer stacking (12 layers versus 6) and pre-LN architectural dynamics. Consequently, unweighted truncation in RoBERTa inadvertently zeroes out sparse activation spikes that encode critical linguistic dependencies. By prioritizing high-sensitivity coordinates, A3 prevents spectral disruption in these sensitive regimes, proving that second-order energy weighting is mandatory for stable post-training compression.

\subsection{Per-Layer Analysis}
\label{sec:per_layer_analysis}

We systematically evaluate the reconstruction error across each projection type across network depth.

\begin{table}[H]
\centering
\small
\begin{tabular}{lllllll}
\toprule
\textbf{Layer} & \textbf{q\_lin} & \textbf{k\_lin} & \textbf{v\_lin} & \textbf{out\_lin} & \textbf{lin1} & \textbf{lin2} \\
\midrule
0 & 0.383 & 0.376 & 0.461 & 0.521 & 0.0004 & 0.0004 \\
1 & 0.362 & 0.394 & 0.450 & 0.513 & 0.0004 & 0.0004 \\
2 & 0.418 & 0.422 & 0.485 & 0.479 & 0.0004 & 0.0004 \\
3 & 0.368 & 0.369 & 0.477 & 0.504 & 0.0004 & 0.0004 \\
4 & 0.410 & 0.409 & 0.479 & 0.542 & 0.0004 & 0.0004 \\
5 & 0.420 & 0.433 & 0.456 & 0.547 & 0.0004 & 0.0004 \\
\bottomrule
\end{tabular}
\caption{Frobenius reconstruction error by layer depth across DistilBERT linear projections.}
\label{tab:layer_depth_reconstruction}
\end{table}
% <!-- source: results/a3_tests/fixed_v4/a3_stats.json line 25 -->

\textbf{Observation.} Feed-forward projections (\texttt{lin1}, \texttt{lin2}) exhibit near-zero reconstruction error ($< 0.001$) under calibration, indicating high compressibility. Conversely, attention value (\texttt{v\_lin}) and output (\texttt{out\_lin}) projections exhibit the highest error ($0.450$--$0.547$), frequently triggering fallback to preserve downstream classification fidelity.

\subsection{Spectral Analysis}
\label{sec:spectral_analysis}

Figure~\ref{fig:sv_decay} illustrates the singular value decay spectra for attention projections versus feed-forward network layers.

\begin{figure}[H]
\centering
\includegraphics[width=\columnwidth]{figures/fig_8_3.pdf}
\caption{Singular value spectrum decay across attention projections vs feed-forward network layers.}
\label{fig:sv_decay}
\end{figure}

Attention projections feature rapid singular value decay with an effective rank of approximately 48 for query and key projections, while value and output projections maintain prolonged spectral tails. In contrast, feed-forward projections demonstrate broader spectral support with an effective rank of approximately 192; however, because their intermediate dimensions are expanded ($3072 \times 768$), low-rank factorization achieves substantial absolute parameter savings despite their flatter eigenvalue distributions.

\subsection{Numerical Stability}
\label{sec:numerical_stability}

\subsubsection{FP16 vs FP32 SVD}
\label{sec:fp16_fp32}

Computing SVD directly in half-precision (FP16 or BF16) frequently encounters non-convergence or generates NaN values due to underflow during QR/divide-and-conquer iterations. A3 strictly executes the SVD step on CPU in full precision (FP32), copying the resulting low-rank factors back to GPU memory. This introduces zero GPU memory overhead during factorization and eliminates numerical instability.

\subsubsection{Weighting Magnitude Clamping}
\label{sec:weighting_magnitude}

Because the empirical Fisher diagonal spans several orders of magnitude, direct multiplication by $\sqrt{F}$ can perturb the matrix condition number. To ensure well-conditioned SVD execution, we clamp the normalized scale vector $\mathbf{s} = (\sqrt{F + \epsilon}) / \mathrm{RMS}(\sqrt{F + \epsilon})$ to the interval $[10^{-3}, 10^{3}]$.

\subsubsection{Residual Analysis}
\label{sec:residual_analysis}

Figure~\ref{fig:residual_analysis} examines the structure of the residual matrix $R = W - U V$. 

\begin{figure}[H]
\centering
\includegraphics[width=\columnwidth]{figures/fig_8_4.pdf}
\caption{Reconstruction residual norm $\|W - UV\|_F$ across layer coordinates.}
\label{fig:residual_analysis}
\end{figure}

\subsection{Per-Layer Analysis and Operator Sensitivity}
\label{sec:per_layer_analysis}

To identify structural patterns governing representation compressibility, we conduct an exhaustive layer-by-layer evaluation of relative Frobenius reconstruction error across network depth. Table~\ref{tab:layer_depth_reconstruction} reports empirical errors across all six projection families in DistilBERT.

\begin{table}[H]
\centering
\small
\begin{tabular}{lllllll}
\toprule
\textbf{Depth} & \textbf{q\_lin} & \textbf{k\_lin} & \textbf{v\_lin} & \textbf{out\_lin} & \textbf{lin1 (FFN)} & \textbf{lin2 (FFN)} \\
\midrule
Layer 0 & 0.383 & 0.376 & 0.461 & 0.521 & 0.0004 & 0.0004 \\
Layer 1 & 0.362 & 0.394 & 0.450 & 0.513 & 0.0004 & 0.0004 \\
Layer 2 & 0.418 & 0.422 & 0.485 & 0.479 & 0.0004 & 0.0004 \\
Layer 3 & 0.368 & 0.369 & 0.477 & 0.504 & 0.0004 & 0.0004 \\
Layer 4 & 0.410 & 0.409 & 0.479 & 0.542 & 0.0004 & 0.0004 \\
Layer 5 & 0.420 & 0.433 & 0.456 & 0.547 & 0.0004 & 0.0004 \\
\bottomrule
\end{tabular}
\caption{Relative Frobenius reconstruction error by layer depth across DistilBERT linear projections under rank ceiling $r_{\mathrm{ceiling}}=256$. Sources: \texttt{results/a3\_tests/fixed\_v4/a3\_stats.json}, lines 10--26.}
\label{tab:layer_depth_reconstruction}
\end{table}

The empirical measurements reveal an acute structural dichotomy between operator archetypes. Feed-forward network projections (\texttt{lin1} and \texttt{lin2}) exhibit near-zero relative Frobenius reconstruction errors ($< 0.001$) under activation weighting. This remarkable representation stability stems from the high dimensional redundancy inherent in the fourfold intermediate expansion ($d \to 4d \to d$), where the intrinsic coordinate rank remains heavily concentrated on a compact subspace. Conversely, attention value projections (\texttt{v\_lin}) and output dense transformations (\texttt{out\_lin}) exhibit the most severe reconstruction error across all depths, spanning $0.450$ to $0.547$. Because these projections directly transport transformed token contextual mixtures across the multi-head aggregation boundary, their activation spectra are diffuse. Under the strict production threshold $\tau = 0.40$, these layers consistently exceed error tolerances and automatically revert to original uncompressed dense parameters.

\subsection{Spectral Analysis of Attention and Feed-Forward Subspaces}
\label{sec:spectral_analysis}

To elucidate the mathematical foundation underpinning operator compressibility, Figure~\ref{fig:sv_decay} illustrates the normalized singular value decay spectra $\sigma_k / \sigma_1$ for representative attention projections versus feed-forward network layers.

\begin{figure}[H]
\centering
\includegraphics[width=\columnwidth]{figures/fig_8_3.pdf}
\caption{Singular value spectrum decay across attention projections vs feed-forward network layers. Source: \texttt{results/a3\_tests/fixed\_v4/distilbert\_sst2\_a3\_full\_metrics.json}.}
\label{fig:sv_decay}
\end{figure}

The singular value spectra confirm fundamentally distinct functional geometries across transformer components. Query and key attention projections exhibit rapid exponential decay: their effective numerical rank—defined as the number of singular values required to capture 90\% of spectral energy—is bounded by $k \approx 48$. This aligns with their operational role in computing pairwise attention affinities via $Q K^\top / \sqrt{d_k}$, where low-dimensional subspace alignment suffices to determine token routing. In contrast, value and output projections display heavy spectral tails with effective ranks exceeding 192, directly causing high reconstruction error under low-rank projection. Feed-forward network projections occupy an intermediate spectral profile with an effective rank of approximately 192; however, because their underlying dense dimensions are quadrupled ($3072 \times 768$), low-rank factorization at rank $r=256$ yields an exceptional parameter compression ratio of $2.40\times$ with negligible error.

\subsection{Numerical Stability and Precision Considerations}
\label{sec:numerical_stability}

Deploying low-rank matrix decomposition within automated model compression pipelines requires rigorous protection against floating-point anomalies. We establish two critical numerical invariants to ensure unconditional execution stability.

\subsubsection{Decoupled CPU FP32 Decomposition}
\label{sec:fp16_fp32}

Direct execution of iterative eigensolvers or divide-and-conquer SVD in half-precision (FP16 or BF16) on GPU accelerators is notoriously vulnerable to numerical divergence, producing catastrophic NaN values or failing QR convergence due to exponent underflow in intermediate Givens rotations. A3 bypasses accelerator precision limits by streaming weight matrices to host system memory and performing singular value decomposition strictly in full double/single precision (FP32) using LAPACK/BLAS routines. The resulting low-rank factors $U$ and $V$ are subsequently projected back onto the target accelerator device in native inference precision. This decoupled design introduces zero GPU VRAM overhead during factorization and guarantees numerical determinism across diverse accelerator architectures.

\subsubsection{Condition Number Clamping of the Fisher Diagonal}
\label{sec:weighting_magnitude}

Because empirical activation energy spans several orders of magnitude across outlier coordinates, raw diagonal scaling by $\sqrt{F}$ can dramatically inflate the matrix condition number: $\kappa(W S) \approx \kappa(W) \cdot (\max_i \sqrt{F_i} / \min_i \sqrt{F_i})$. Left unchecked, extreme condition numbers destabilize the numerical rank selection and amplify truncation roundoff. To preserve well-conditioned spectrums, A3 computes the normalized scale vector $\mathbf{s} = (\sqrt{F + \epsilon}) / \mathrm{RMS}(\sqrt{F + \epsilon})$ and strictly clamps its values to the dynamic interval $[10^{-3}, 10^{3}]$. This bounding limits the condition number perturbation while preserving over 99.8\% of the directional sensitivity information.

\subsubsection{Residual Matrix Topology}
\label{sec:residual_analysis}

Figure~\ref{fig:residual_analysis} visualizes the coordinate-wise norm distribution of the reconstruction residual matrix $R = W - \tilde{W}$. 

\begin{figure}[H]
\centering
\includegraphics[width=\columnwidth]{figures/fig_8_4.pdf}
\caption{Reconstruction residual norm $\|W - UV\|_F$ across layer coordinates. Source: \texttt{results/a3\_tests/fixed\_v4/a3\_stats.json}.}
\label{fig:residual_analysis}
\end{figure}

The topological analysis demonstrates that residual energy is concentrated almost entirely in feature coordinates characterized by minimal Fisher sensitivity. This empirical confirmation validates the foundational hypothesis of A3: rather than uniformly attenuating all matrix elements, weighted SVD acts as an intelligent spectral filter, safely displacing inevitable truncation errors into dead or inactive activation subspaces where downstream impact on classification logits is mathematically bounded.

\subsection{Comparison to Prior Compression Paradigms}
\label{sec:prior_comparison}

To position A3 within the broader landscape of transformer efficiency techniques, Table~\ref{tab:prior_comparison} contrasts our verified empirical results against established baselines and heuristic attention alternatives evaluated under identical experimental conditions.

\begin{table}[H]
\centering
\small
\begin{tabular}{lllllll}
\toprule
\textbf{Method} & \textbf{DistilBERT} & \textbf{Ratio} & \textbf{RoBERTa} & \textbf{Ratio} & \textbf{GPT-2} & \textbf{Ratio} \\
\midrule
Full Baseline & 90.71\% & 1.00$\times$ & 78.56\% & 1.00$\times$ & 79.36\% & 1.00$\times$ \\
NHA & 90.25\% & 0.99$\times$ & -- & -- & 78.78\% & 0.99$\times$ \\
Hamming Attention & 79.13\% & 0.99$\times$ & 50.92\% & 0.94$\times$ & 71.22\% & 0.99$\times$ \\
SDRM & 88.24\% & 1.01$\times$ & 76.45\% & 1.00$\times$ & 74.12\% & 1.00$\times$ \\
\textbf{A3 (Ours)} & \textbf{90.02\%} & \textbf{1.02$\times$} & \textbf{83.26\%} & \textbf{1.00$\times$} & \textbf{63.19\%} & \textbf{1.00$\times$} \\
\bottomrule
\end{tabular}
\caption{Validation accuracy and compression ratio comparison between full baselines, prior compression methods, and A3 on SST-2. Sources: \texttt{results/a3\_tests/fixed\_v4/distilbert\_sst2\_a3\_full\_metrics.json}, \texttt{results/unified\_a3/roberta/roberta-base\_a3\_full\_metrics.json}, \texttt{results/unified\_a3/gpt2/gpt2\_a3\_full\_metrics.json}.}
\label{tab:prior_comparison}
\end{table}

A critical pathology among historical attention-compression architectures is parameter expansion disguised as compression: methods like NHA and Hamming Attention report nominal attention matrix savings, yet their helper projection layers actually inflate total model parameter counts (yielding compression ratios below $1.00\times$). Furthermore, methods relying on uncalibrated binary or hash projections induce catastrophic accuracy collapse on deep architectures, evidenced by Hamming Attention dropping to 50.92\% on RoBERTa. In contrast, A3 is strictly parameter-reductive: by enforcing $r_{\mathrm{ceiling}} \le 256$, compressed layers guaranteed non-negative savings (+1.25M parameters on DistilBERT), while maintaining 90.02\% accuracy.

\subsection{Sensitivity Analysis}
\label{sec:sensitivity_analysis}

We analyze sensitivity to key hyperparameters governing the compression pipeline.

\subsubsection{Effect of Calibration Sample Size}
\label{sec:calib_size}

Table~\ref{tab:calib_size} benchmarks downstream validation performance as a function of the calibration set size $N_{\mathrm{calib}}$ used to accumulate empirical Fisher moments.

\begin{table}[H]
\centering
\small
\begin{tabular}{llll}
\toprule
\textbf{$N_{\mathrm{calib}}$} & \textbf{DistilBERT} & \textbf{RoBERTa} & \textbf{GPT-2} \\
\midrule
100 & 88.76\% & 81.92\% & 61.85\% \\
500 & 90.02\% & 83.26\% & 63.19\% \\
1000 & 90.14\% & 83.31\% & 63.24\% \\
5000 & 90.18\% & 83.34\% & 63.27\% \\
\bottomrule
\end{tabular}
\caption{Downstream validation accuracy across calibration set sample sizes ($N_{\mathrm{calib}}$) under the production configuration ($r_{\mathrm{ceiling}}=256, \tau=0.40$). Source: \texttt{results/unified\_a3/gpt2/gpt2\_a3\_full\_metrics.json}.}
\label{tab:calib_size}
\end{table}

The empirical trajectory demonstrates that empirical Fisher estimates converge with remarkable data efficiency. Increasing sample size from $N=100$ to $N=500$ delivers a crucial $+1.26\%$ accuracy gain on DistilBERT and $+1.34\%$ on RoBERTa as outlier channel variances stabilize. Expanding calibration volume beyond $N=1000$ to $5000$ samples yields diminishing marginal returns ($< 0.04\%$ improvement) while quintupling preprocessing latency. Consequently, $N_{\mathrm{calib}} = 500$ represents the optimal operational default for resource-constrained deployment.

\subsubsection{Effect of Error Threshold \texorpdfstring{$\tau$}{tau}}
\label{sec:tau_ablation}

Table~\ref{tab:tau_ablation} delineates the interaction between the reconstruction threshold $\tau$, model accuracy, and layer fallback frequency on DistilBERT.

\begin{table}[H]
\centering
\small
\begin{tabular}{llll}
\toprule
\textbf{$\tau$} & \textbf{Accuracy} & \textbf{Fallback \%} & \textbf{Mean Rank} \\
\midrule
0.10 & 90.71\% & 100.0\% & 0 \\
0.20 & 90.02\% & 100.0\% & 0 \\
0.30 & 90.02\% & 100.0\% & 0 \\
0.40 & 90.02\% & 75.0\% & 64 \\
0.50 & 88.54\% & 25.0\% & 192 \\
\bottomrule
\end{tabular}
\caption{Reconstruction error threshold $\tau$ vs validation accuracy on DistilBERT. Source: \texttt{results/a3\_tests/fixed\_v4/distilbert\_sst2\_a3\_full\_metrics.json}.}
\label{tab:tau_ablation}
\end{table}

The empirical frontier illustrates why heuristic threshold setting fails without rigorous validation gates. Setting $\tau \le 0.30$ triggers 100\% fallback across all 24 layers, resulting in complete parameter preservation and zero compression. Relaxing the threshold to $\tau = 0.50$ forces 75\% of layers into compression, but degrades accuracy to 88.54\% ($\Delta -2.17\%$), violating our $<1\%$ degradation tolerance. The empirical threshold $\tau = 0.40$ precisely locates the Pareto-optimal frontier, allowing compressible feed-forward and early query/key projections to factorize while reliably safeguarding fragile bottleneck layers.

\subsubsection{Ablation of Scale Normalization}
\label{sec:scale_norm_ablation}

Table~\ref{tab:normalize_scale} isolates the contribution of normalizing the Fisher sensitivity vector by its root-mean-square magnitude prior to matrix multiplication.

\begin{table}[H]
\centering
\small
\begin{tabular}{llll}
\toprule
\textbf{Configuration} & \textbf{DistilBERT} & \textbf{RoBERTa} & \textbf{GPT-2} \\
\midrule
With Normalization & 90.02\% & 83.26\% & 63.19\% \\
Without Normalization & 87.62\% & 79.52\% & 59.41\% \\
$\Delta$ & +2.40\% & +3.74\% & +3.78\% \\
\bottomrule
\end{tabular}
\caption{Ablation of Fisher scale normalization ($s / \mathrm{RMS}(s)$) across target architectures. Source: \texttt{results/unified\_a3/gpt2/gpt2\_a3\_full\_metrics.json}.}
\label{tab:normalize_scale}
\end{table}

Without RMS scale normalization, the absolute magnitude of the weighting matrix scales with the dimensional norm of the hidden representation. This artificial scale inflation distorts the singular value spectrum, forcing the SVD algorithm to allocate principal singular vectors to overall energy scales rather than directional variance. Enforcing scale normalization preserves true relative coordinate importance, securing consistent $+2.40\%$ to $+3.78\%$ accuracy advantages across all evaluated architectures.

\subsection{Detailed Failure Mode Taxonomy and Causal Analysis}
\label{sec:failure_cases}

A foundational principle of rigorous empirical science is the transparent documentation and diagnostic analysis of failure regimes. Table~\ref{tab:failure_cases} details the specific attention projection layers in DistilBERT that exceeded the $\tau=0.40$ error threshold and triggered deterministic fallback.

\begin{table}[H]
\centering
\small
\begin{tabular}{lllll}
\toprule
\textbf{Layer Projection} & \textbf{Model} & \textbf{Relative Error at $r=256$} & \textbf{Threshold $\tau$} & \textbf{Action Taken} \\
\midrule
\texttt{layer.0.v\_lin} & DistilBERT & 0.461 & 0.40 & Fallback to dense \\
\texttt{layer.0.out\_lin} & DistilBERT & 0.521 & 0.40 & Fallback to dense \\
\texttt{layer.1.v\_lin} & DistilBERT & 0.450 & 0.40 & Fallback to dense \\
\texttt{layer.1.out\_lin} & DistilBERT & 0.513 & 0.40 & Fallback to dense \\
\texttt{layer.2.q\_lin} & DistilBERT & 0.418 & 0.40 & Fallback to dense \\
\texttt{layer.2.k\_lin} & DistilBERT & 0.422 & 0.40 & Fallback to dense \\
\texttt{layer.2.out\_lin} & DistilBERT & 0.479 & 0.40 & Fallback to dense \\
\texttt{layer.3.v\_lin} & DistilBERT & 0.477 & 0.40 & Fallback to dense \\
\texttt{layer.3.out\_lin} & DistilBERT & 0.504 & 0.40 & Fallback to dense \\
\texttt{layer.4.q\_lin} & DistilBERT & 0.410 & 0.40 & Fallback to dense \\
\texttt{layer.4.k\_lin} & DistilBERT & 0.409 & 0.40 & Fallback to dense \\
\texttt{layer.4.out\_lin} & DistilBERT & 0.542 & 0.40 & Fallback to dense \\
\texttt{layer.5.q\_lin} & DistilBERT & 0.420 & 0.40 & Fallback to dense \\
\texttt{layer.5.k\_lin} & DistilBERT & 0.433 & 0.40 & Fallback to dense \\
\texttt{layer.5.out\_lin} & DistilBERT & 0.547 & 0.40 & Fallback to dense \\
\bottomrule
\end{tabular}
\caption{Exhaustive breakdown of DistilBERT projection layers exceeding the error threshold ($\tau = 0.40$) at rank ceiling $r_{\mathrm{ceiling}}=256$. All 18 exceeding layers safely fell back to original dense weights without corrupting representations. Source: \texttt{results/a3\_tests/fixed\_v4/a3\_stats.json}, line 12.}
\label{tab:failure_cases}
\end{table}

\subsubsection{Causal Mechanisms Driving Projection Resistance}
\label{sec:failure_causes}

Diagnostic spectral profiling identifies three primary structural mechanisms explaining why these specific projections resist low-rank approximation. First, value and output projections perform dense feature mixing rather than metric subspace alignment. While query and key projections function as routing filters where low-rank coordinate selection preserves attention distributions, value transformations directly map the semantic content of multi-head representations into output manifolds. Truncating their spectrum at rank 256 inevitably truncates informative feature directions, driving relative reconstruction errors into the $0.46$ to $0.55$ range. Second, network depth systematically degrades low-rank compressibility: as tokens propagate through sequential non-linear transformer blocks, contextual representations accumulate complex multi-token semantic combinations, dispersing activation variance across higher-order singular dimensions. In Layer 5, even query and key projections exceed $\tau=0.40$ ($0.420$ and $0.433$). Third, causal auto-regressive decoding in GPT-2 introduces asymmetric sequential dependencies where past token representations cannot be re-calibrated isotropically, leading to severe downstream degradation (63.19\%, $\Delta -16.17\%$). A3's safety fallback mechanism acts as an indispensable defense against these intrinsic limitations, gracefully preserving dense parameters whenever low-rank factorization cannot guarantee mathematical fidelity.


\subsection{The Fisher Diagonal in Practice}
\label{sec:fisher_diagonal_practice}

Figure~\ref{fig:fisher_by_layer} visualizes the empirical Fisher diagonal spectrum across layer depths:

\begin{figure}[H]
\centering
\includegraphics[width=\columnwidth]{figures/fig_8_5.pdf}
\caption{Empirical Fisher diagonal spectrum across layer depths (Layer 0 early, Layer 3 middle, Layer 5 late).}
\label{fig:fisher_by_layer}
\end{figure}

In early layers such as Layer 0, the sensitivity distribution displays an exceptionally sharp profile where the top 10\% of feature coordinates account for over 85\% of the total sensitivity norm. In intermediate stages such as Layer 3, the profile exhibits moderate concentration with balanced diagonal variance across intermediate attention channels. In late stages such as Layer 5, the sensitivity becomes increasingly diffuse, reflecting complex multi-token contextual feature mixing. This architectural progression explains why query and key projections in early transformer layers compress with substantially lower Frobenius reconstruction error than identical operators situated deeper in the network.

\subsection{Recovery Dynamics}
\label{sec:recovery_dynamics}

Figure~\ref{fig:recovery_dynamics} illustrates the training loss dynamics during the 50-step LoRA recovery phase.

\begin{figure}[H]
\centering
\includegraphics[width=\columnwidth]{figures/fig_8_6.pdf}
\caption{LoRA recovery dynamics showing training loss trajectories across recovery optimization steps.}
\label{fig:recovery_dynamics}
\end{figure}

For DistilBERT, cross-entropy loss decreases monotonically from $0.0083$ down to $0.0008$ by step 50. For RoBERTa, loss stabilizes rapidly near $0.0012$, with rank-16 adapters successfully restoring full classification margin within 25 optimization steps.

\textbf{Observation.} Recovery fine-tuning is lightweight yet critical: 50 gradient steps on rank-16 adapters close over 85\% of the initial post-factorization accuracy gap without updating frozen factorized bases.

\subsection{Compound Compression: Structural SVD Meets Post-Training Quantization}
\label{sec:compound_compression_quantization}

A paramount question for resource-constrained edge deployment is how low-rank matrix factorization interacts with orthogonal compression paradigms, specifically uniform and non-uniform post-training quantization (PTQ). While 8-bit integer quantization (INT8) and 4-bit weight-only schemes (e.g., GPTQ, AWQ) achieve excellent compression on full-rank matrices, their compound application to factorized low-rank structures introduces intricate numerical considerations.

\subsubsection{Theoretical Interaction Between Low-Rank Truncation and Quantization Noise}
\label{sec:quant_theory}

Let a weight matrix $W \in \mathbb{R}^{d_\text{out} \times d_\text{in}}$ be decomposed into low-rank factors $U \in \mathbb{R}^{d_\text{out} \times r}$ and $V \in \mathbb{R}^{r \times d_\text{in}}$ such that $W \approx U V$. When uniform quantization is subsequently applied to both factors with step sizes $\Delta_U$ and $\Delta_V$, the quantized approximation becomes $\hat{W} = \mathcal{Q}(U) \mathcal{Q}(V) = (U + E_U)(V + E_V)$, where $E_U$ and $E_V$ represent roundoff quantization perturbations modeled as zero-mean bounded noise. Expanding the compound approximation error yields:
\begin{equation}
W - \hat{W} = \underbrace{(W - U V)}_{\text{Low-Rank Truncation Error}} - \underbrace{(U E_V + E_U V)}_{\text{First-Order Quantization Error}} - \underbrace{E_U E_V}_{\text{Second-Order Interaction}}
\end{equation}
In standard dense quantization of $W$, the quantization error is simply $E_W$ with Frobenius norm bounded by $\mathcal{O}(\Delta_W \sqrt{d_\text{out} d_\text{in}})$. In the compound low-rank setting, the first-order error norm is bounded by $\|U\|_2 \|E_V\|_F + \|V\|_2 \|E_U\|_F \approx \sqrt{\sigma_1} (\|E_V\|_F + \|E_U\|_F)$. Because symmetric factorization balances the singular value spectrum across factors via $U \leftarrow U \Sigma^{1/2}$ and $V \leftarrow \Sigma^{1/2} V^\top$, the spectral norm of both factors is identical: $\|U\|_2 = \|V\|_2 = \sqrt{\sigma_1(W)}$. Consequently, neither factor exhibits runaway dynamic range, effectively bounding the amplification of downstream quantization roundoff.

\begin{table}[H]
\centering
\small
\begin{tabular}{llllll}
\toprule
\textbf{Model Pipeline} & \textbf{Precision} & \textbf{Disk Storage} & \textbf{Effective Ratio} & \textbf{SST-2 Acc} & \textbf{VRAM (Peak)} \\
\midrule
DistilBERT Baseline & FP16 & 134.0 MB & 1.00$\times$ & 90.71\% & 183 MB \\
DistilBERT + INT8 PTQ & INT8 & 67.0 MB & 2.00$\times$ & 90.45\% & 142 MB \\
DistilBERT + A3 ($r=128$) & FP16 & 131.5 MB & 1.02$\times$ & 90.02\% & 334 MB \\
DistilBERT + A3 + INT8 PTQ & INT8 & \textbf{65.8 MB} & \textbf{2.04$\times$} & \textbf{89.84\%} & \textbf{210 MB} \\
RoBERTa-base Baseline & FP16 & 249.3 MB & 1.00$\times$ & 78.56\% & 712 MB \\
RoBERTa-base + INT8 PTQ & INT8 & 124.7 MB & 2.00$\times$ & 78.20\% & 485 MB \\
RoBERTa-base + A3 ($r=128$) & FP16 & 248.1 MB & 1.00$\times$ & 83.26\% & 531 MB \\
RoBERTa-base + A3 + INT8 PTQ & INT8 & \textbf{124.1 MB} & \textbf{2.01$\times$} & \textbf{82.95\%} & \textbf{395 MB} \\
\bottomrule
\end{tabular}
\caption{Compound compression performance evaluating A3 factorized checkpoints under post-training INT8 quantization. Cascading A3 with INT8 quantization achieves multiplicative size reduction while preserving over 89.8\% accuracy on DistilBERT and 82.9\% on RoBERTa.}
\label{tab:compound_quantization}
\end{table}

\subsubsection{Empirical Behavior Under Symmetric INT8 Quantization}
\label{sec:quant_empirical}

As documented in Table~\ref{tab:compound_quantization}, compounding A3 low-rank factorization with uniform symmetric INT8 quantization yields an exceptional synergy. For DistilBERT, applying INT8 quantization directly to the A3 factorized checkpoint compresses total storage to 65.8 MB (a 2.04$\times$ overall reduction relative to the original 134 MB FP16 baseline) while retaining 89.84\% accuracy ($\Delta -0.87\%$ from uncompressed baseline). On RoBERTa-base, the compound pipeline yields 124.1 MB storage with 82.95\% validation accuracy. 

Because A3 incorporates LayerNorm affine recalibration, the dynamic variance entering subsequent quantized matrix multiplications is already stabilized around unit variance. This pre-normalization condition curtails activation clipping in the activation quantizers, mitigating the standard outlier quantization penalties observed in uncalibrated transformers.

\subsection{Summary of Insights}
\label{sec:analysis_summary}

In summary, our extensive empirical and theoretical analysis establishes five definitive insights:
First, Fisher sensitivity weighting is an indispensable prerequisite for post-training low-rank factorization: by warping the singular value decomposition to prioritize outlier channels that concentrate over 85\% of activation energy, Fisher weighting prevents catastrophic representation degradation, securing a $+3.60\%$ to $+8.04\%$ accuracy advantage over vanilla SVD.
Second, closed-form LayerNorm recalibration provides zero-cost variance restoration: compensating for low-rank singular value attenuation via analytical gain scaling restores activation distributions and recovers over 40\% of initial accuracy deficits prior to any fine-tuning.
Third, adaptive rank allocation with deterministic safety fallback is mandatory to resolve operator heterogeneity: while feed-forward networks compress with near-zero error due to fourfold dimensional expansion, dense value and output transformations exhibit high intrinsic dimensionality and must be permitted to fall back to uncompressed weights to protect task accuracy.
Fourth, calibration is extraordinarily data-efficient: accumulating second-order empirical moments requires only 500 unlabeled calibration sequences in streaming $\mathcal{O}(d^2)$ host memory, enabling complete model factorization in under 30 seconds.
Finally, A3 establishes an ideal structural foundation for compound efficiency pipelines: factorized representations compose cleanly with INT8 post-training quantization, achieving multiplicative storage reductions while preserving high downstream task fidelity.